 \documentclass[french]{book}
\usepackage[utf8]{inputenc}
\usepackage[french]{babel}
%\usepackage{teach}
\usepackage{verbatim}
\usepackage{amsmath,amsfonts,amssymb}
\usepackage{pgf,pgfplots,tikz,tkz-tab}
\usepackage{mathrsfs}
\usepackage{multicol}
\usepackage{lscape}
\usepackage{pstricks-add}
%\pgfplotsset{compat=1.15}
\usetikzlibrary{arrows}
%\redefinecolor{thm}{title}{bgcolor}{alizarin}
%\redefinecolor{prop}{title}{bgcolor}{meatbrown}
%\redefinecolor{defn}{title}{bgcolor}{olivine}
\definecolor{alizarin}{rgb}{0.82, 0.1, 0.26}
\definecolor{meatbrown}{rgb}{0.9, 0.72, 0.23}
\definecolor{olivine}{rgb}{0.6, 0.73, 0.45}
\usepackage{geometry}
%\geometry{top=1cm, bottom=1cm, left=2cm, right=2cm}
\geometry{hmargin=2cm,vmargin=2cm}
\pagestyle{empty}

\begin{document}
\section*{Devoir à la maison II}
\begin{center}
\definecolor{uuuuuu}{rgb}{0.26666666666666666,0.26666666666666666,0.26666666666666666}
\definecolor{ududff}{rgb}{0.30196078431372547,0.30196078431372547,1}
\begin{tikzpicture}[scale=0.7,line cap=round,line join=round,>=triangle 45,x=1cm,y=1cm]
\clip(-12.248318523626864,-6.148614831021472) rectangle (10.42454346131598,6.622583208934634);
\draw [line width=1pt] (-3.62,5.03)-- (-5.02,-1.97);
\draw [line width=1pt] (-3.62,5.03)-- (1.8159411764079942,2.1057239853717946);
\draw [line width=1pt] (-4.32,1.53)-- (-0.9020294117960029,3.5678619926858977);
\draw [line width=1pt] (-5.02,-1.97)-- (1.8159411764079942,2.1057239853717946);
\begin{scriptsize}
\draw [fill=ududff] (-3.62,5.03) circle (2.5pt);
\draw[color=ududff] (-3.4802976778872483,5.409230829271675) node {$A$};
\draw [fill=ududff] (-5.02,-1.97) circle (2.5pt);
\draw[color=ududff] (-5.2,-1.5874726738942917) node {$B$};
\draw [fill=ududff] (1.8159411764079942,2.1057239853717946) circle (2.5pt);
\draw[color=ududff] (1.9576465638138871,2.486557214025132) node {$C$};
\draw [fill=uuuuuu] (-4.32,1.53) circle (2pt);
\draw[color=uuuuuu] (-4.471111441490976,1.8665961441243502) node {$I$};
\draw [fill=uuuuuu] (-0.9020294117960029,3.5678619926858977) circle (2pt);
\draw[color=uuuuuu] (-0.7524689703238123,3.921324261509799) node {$J$};
\end{scriptsize}
\end{tikzpicture}

\end{center}
\subsection*{Exercice 1 - Le théorème des milieux}
Soient ABC un triangle quelconque; $I$,$J$ les milieux respectifs [AB] et [BC].

\underline{L'objectif ici est de montrer que $\overrightarrow{IJ}$ est colinéaire à $\overrightarrow{BC}$.}\\


\begin{itemize}
\item[A)] Exprimer $\overrightarrow{BA}$ en fonction de $\overrightarrow{IA}$.
\item[B)] Exprimer $\overrightarrow{AC}$ en fonction de $\overrightarrow{AJ}$.
\item[C)]A l'aide de la relation de Chasles et \underline{des points $A,B$ et $C$} recopier et compléter l'expression ci dessous:  $$\overrightarrow{BC} = \overrightarrow{\cdots} + \overrightarrow{\cdots}$$
\item[D)] A l'aide de C) puis de A) et B) ainsi que la relation de Chasles, montrer que $\overrightarrow{BC}=2\overrightarrow{IJ}$.
\item[E)] Conclure.
\end{itemize}

\subsection*{Exercice 2 - Le parralèlogramme de Varignon ( \underline{bonus} )}

\underline{Dans cet exercice, on supose que le résultat de l'exercice 1 est vrai, c'est à dire $2\overrightarrow{IJ}=\overrightarrow{BC}$.}\\


\begin{itemize}
\item[A)] Tracer le quadrilatère ABDC de votre choix ( noté dans le sens inverse des aiguilles d'une montre  ).\\
\end{itemize}
On note I,J,K,L les milieux respectifs de [AB], [AC], [BD], [DC].\\


\begin{itemize}
\item[B)] Tracer I,J,K,L au compas, laissez les traits de construction.\\

\end{itemize}

\underline{L'objectif est de montrer que IJKL est un parralèlogramme.}\\



\begin{itemize}
\item[C)] Expliquer \underline{à l'aide de l'exercice 1} que $\overrightarrow{KL}$ est colinéaire à  $\overrightarrow{BC}$. 
\item[D)] Exprimer $\overrightarrow{BC}$ en fonction de $\overrightarrow{KL}$.
\item[E)] En déduire que  $\overrightarrow{IJ}=\overrightarrow{KL}$.
\item[F)] Conclure.
\end{itemize}

\end{document}